% Volume IV: Optimization and Computation
% Continuity note for Chapter 22.
% Previous dependencies: Chapter 20 supplies objective, feasible-set, convexity, subgradient, and regularization geometry; Chapter 21 supplies gradient and Newton update intuition; Chapters 7-10 supply gradients, Hessians, differentials, Taylor approximation, and local linear maps; Chapters 15 and 18 supply expectation, Bayes, entropy, KL divergence, and cross-entropy for later variational inference.
% New durable concepts: equality constraints, tangent spaces, Lagrange multipliers, Lagrangian functions, active and inactive inequalities, KKT stationarity, primal feasibility, dual feasibility, complementary slackness, constraint qualifications, dual functions, weak and strong duality, Slater intuition, shadow prices, functionals, variations, first variation, functional derivatives, fundamental lemma, Euler-Lagrange equations, constrained variational problems, ELBO stationarity, optimal-control and neural-field energy links.
% Forward hooks: Variational inference, ELBOs, energy-based models, optimal control, neural fields, predictive coding, active inference, and statistical-physics free-energy ideas all reuse Lagrangians, dual bounds, constrained stationarity, or functional derivatives.
% Notation commitments: Finite-dimensional variables are theta in R^d; equality constraints are h(theta)=0; inequality constraints are g_i(theta)<=0; Lagrange multipliers for inequalities are lambda_i>=0 and for equalities are nu_j; Lagrangians are mathcal L; function-valued decision variables are y or q; variations are eta; functional derivatives are delta J/delta y.
% Figure plan: no figures are included in this authored source. Proposed raster figures are listed in imagegen-figures/ch22-figure-metadata.md.
\section{Chapter 22: Constrained Optimization, Duality, and Variational Principles}
\label{ch:constrained-optimization-duality-variational-principles}

Many scientific optimization problems are not free searches over all vectors. Probabilities must sum to one. Firing rates may be nonnegative. Control policies may obey resource limits. Neural fields and physical systems often minimize energies over entire functions, not just finite-dimensional vectors. Variational inference optimizes over distributions while maintaining normalization.

This chapter extends optimization from Chapter 20 in two directions. First, it adds constraints and the multipliers that measure the cost of satisfying them. Second, it replaces finite-dimensional vectors by functions and distributions, leading to variational calculus. The same idea runs throughout: an optimum is a point where all feasible first-order changes fail to decrease the objective.

\paragraph{Chapter problem.}
How do constraints, dual variables, and variations turn first-order optimality into usable equations for finite-dimensional problems, probability distributions, controls, energy models, and neural fields?

\paragraph{Section map.}
Section 22.1 develops Lagrange multipliers for equality constraints. Section 22.2 adds inequality constraints and the KKT conditions. Section 22.3 derives dual functions, weak duality, strong-duality intuition, and shadow prices. Section 22.4 develops variational calculus, functional derivatives, Euler-Lagrange equations, and links to ELBOs, optimal control, free energy, and neural fields.

\subsection{22.1 Lagrange multipliers}

\paragraph{Guiding question.}
When an optimum must stay on a constraint surface, how can gradients reveal the best feasible point?

\paragraph{Beginner-facing intuition.}
In unconstrained optimization, a smooth local minimum has zero gradient because every small direction is feasible. With an equality constraint, most directions are illegal. At a constrained optimum, the objective need not be flat in every direction; it only needs to be flat along directions tangent to the constraint surface.

The Lagrange multiplier condition says that the objective gradient must be built from the constraint gradients. Geometrically, the objective's level surface is tangent to the constraint surface at the optimum.

\paragraph{Formal definitions and notation.}
Consider
\[
\min_{\theta\in\mathbb R^d} f(\theta)
\quad\text{subject to}\quad
h(\theta)=0,
\]
where \(h:\mathbb R^d\to\mathbb R^m\) and \(f\) are differentiable. The feasible set is
\[
C=\{\theta:h(\theta)=0\}.
\]
The Jacobian \(J_h(\theta)\in\mathbb R^{m\times d}\) has rows \(\nabla h_j(\theta)^\top\).

A feasible perturbation \(v\) tangent to the constraint at a regular point satisfies
\[
J_h(\theta)v=0.
\]
The Lagrangian is
\[
\boxed{
\mathcal L(\theta,\nu)=f(\theta)+\nu^\top h(\theta),
}
\]
where \(\nu\in\mathbb R^m\) is a vector of Lagrange multipliers. Some texts use a minus sign. The sign convention changes \(\nu\), not the content.

\paragraph{Core mechanism: constrained first-order stationarity.}
Assume \(\theta^\star\) is a local constrained minimizer and the constraint gradients are linearly independent at \(\theta^\star\). For every tangent direction \(v\) with \(J_h(\theta^\star)v=0\), the directional derivative of \(f\) must vanish:
\[
\nabla f(\theta^\star)^\top v=0.
\]
Thus \(\nabla f(\theta^\star)\) is orthogonal to the tangent space. Under the linear-independence condition, the orthogonal complement of the tangent space is the span of the constraint gradients. Therefore there exists \(\nu^\star\in\mathbb R^m\) such that
\[
\boxed{
\nabla f(\theta^\star)+J_h(\theta^\star)^\top\nu^\star=0,
\qquad
h(\theta^\star)=0.
}
\]
Equivalently,
\[
\nabla_\theta\mathcal L(\theta^\star,\nu^\star)=0,
\qquad
\nabla_\nu\mathcal L(\theta^\star,\nu^\star)=h(\theta^\star)=0.
\]
The multiplier \(\nu\) measures how much the optimum value would locally change if the equality constraint were relaxed or tightened.

\paragraph{Concrete worked example.}
Find the point on the unit circle closest to \(a=(2,1)\). Minimize
\[
f(x,y)=\frac12\bigl((x-2)^2+(y-1)^2\bigr)
\]
subject to
\[
h(x,y)=x^2+y^2-1=0.
\]
The Lagrangian is
\[
\mathcal L(x,y,\nu)=
\frac12\bigl((x-2)^2+(y-1)^2\bigr)+\nu(x^2+y^2-1).
\]
Stationarity gives
\[
x-2+2\nu x=0,\qquad y-1+2\nu y=0,
\]
so
\[
(1+2\nu)(x,y)=(2,1).
\]
Thus \((x,y)\) is a positive scalar multiple of \((2,1)\). The unit-circle constraint gives
\[
(x,y)=\frac{(2,1)}{\sqrt{5}}.
\]
The gradient of the objective and the gradient of the constraint are aligned at the closest point.

\paragraph{Computational neuroscience example.}
Suppose a neural population model estimates nonnegative firing probabilities \(p_1,\ldots,p_K\) over \(K\) latent states and requires
\[
\sum_{k=1}^K p_k=1.
\]
If the objective is a smooth negative log-likelihood plus penalties, a Lagrange multiplier enforces normalization. The stationarity condition says that, at the optimum, the unconstrained gradient is balanced by the same multiplier in every probability coordinate, after accounting for any active nonnegativity constraints handled by KKT conditions.

\paragraph{AI and machine-learning example.}
Softmax probabilities solve a constrained normalization problem. If one maximizes
\[
\sum_k p_k a_k+\tau H(p)
\]
over \(p_k\ge 0\) and \(\sum_k p_k=1\), the equality constraint multiplier produces the normalization constant. The solution has the form
\[
p_k=\frac{\exp(a_k/\tau)}{\sum_j\exp(a_j/\tau)}.
\]
This connects constrained optimization to probabilistic choice, attention weights, and entropy-regularized policies.

\paragraph{Common misconceptions and failure modes.}
\begin{itemize}[leftmargin=*]
\item The constrained optimum need not satisfy \(\nabla f=0\); only feasible first-order directions must fail to decrease \(f\).
\item Lagrange multiplier equations are necessary conditions, not automatic guarantees of a minimum.
\item Constraint qualifications matter. If constraint gradients are dependent or zero, multiplier conclusions can fail.
\item Multipliers depend on the scaling and sign convention of the constraint equation.
\item Equality constraints and inequality constraints have different multiplier rules.
\end{itemize}

\paragraph{Exercises.}
\begin{enumerate}[label=\textbf{22.1.\arabic*.}, leftmargin=*]
\item \textbf{Mechanical.} Write the Lagrangian for minimizing \(x^2+y^2\) subject to \(x+y=1\).
\item \textbf{Proof or derivation.} Derive the condition \(\nabla f(\theta^\star)+J_h(\theta^\star)^\top\nu^\star=0\) from feasible tangent directions.
\item \textbf{Numerical/computational.} Solve \(\min_{x,y}x^2+y^2\) subject to \(x+y=1\) using Lagrange multipliers.
\item \textbf{Interpretation.} In a normalized probability model, what does the equality constraint \(\sum_k p_k=1\) prevent?
\item \textbf{Misconception/debugging.} A student finds \(\nabla f\neq 0\) at a constrained optimum and says the solution must be wrong. Explain why this can be correct.
\end{enumerate}

\subsection{22.2 KKT conditions}

\paragraph{Guiding question.}
How do first-order conditions change when constraints can be active or inactive inequalities?

\paragraph{Beginner-facing intuition.}
An inequality constraint such as \(g(\theta)\le 0\) may or may not matter at the optimum. If the solution lies strictly inside the feasible region, the constraint is inactive and should exert no force. If the solution lies on the boundary, the constraint can push back against the objective gradient.

The Karush-Kuhn-Tucker, or KKT, conditions encode this logic. They combine stationarity, feasibility, nonnegative multipliers for inequalities, and complementary slackness: an inequality multiplier is positive only when its constraint is tight.

\paragraph{Formal definitions and notation.}
Consider
\[
\min_{\theta\in\mathbb R^d} f(\theta)
\]
subject to
\[
g_i(\theta)\le 0,\qquad i=1,\ldots,m,
\]
\[
h_j(\theta)=0,\qquad j=1,\ldots,r.
\]
The Lagrangian is
\[
\boxed{
\mathcal L(\theta,\lambda,\nu)
=f(\theta)+\sum_{i=1}^m\lambda_i g_i(\theta)
+\sum_{j=1}^r\nu_j h_j(\theta),
}
\]
where inequality multipliers satisfy \(\lambda_i\ge 0\), while equality multipliers \(\nu_j\) are unrestricted.

Under appropriate constraint qualifications, a local optimum \(\theta^\star\) has multipliers \(\lambda^\star,\nu^\star\) satisfying:
\[
\text{Stationarity:}\quad
\nabla_\theta \mathcal L(\theta^\star,\lambda^\star,\nu^\star)=0.
\]
\[
\text{Primal feasibility:}\quad
g_i(\theta^\star)\le 0,\quad h_j(\theta^\star)=0.
\]
\[
\text{Dual feasibility:}\quad
\lambda_i^\star\ge 0.
\]
\[
\text{Complementary slackness:}\quad
\lambda_i^\star g_i(\theta^\star)=0
\quad\text{for each }i.
\]

\paragraph{Core mechanism: active constraints.}
An inequality is active at \(\theta^\star\) if
\[
g_i(\theta^\star)=0.
\]
It is inactive if
\[
g_i(\theta^\star)<0.
\]
Complementary slackness says:
\[
g_i(\theta^\star)<0\quad\Longrightarrow\quad \lambda_i^\star=0.
\]
Only active constraints can have nonzero multipliers. At an active boundary, \(\lambda_i^\star\nabla g_i(\theta^\star)\) contributes a normal force that balances the objective gradient.

For convex problems, KKT conditions become especially powerful. If \(f\) and the \(g_i\) are convex, the \(h_j\) are affine, and a regularity condition such as Slater's condition holds, then any point satisfying KKT is a global optimum.

\paragraph{Concrete worked example.}
Solve
\[
\min_x (x-2)^2
\quad\text{subject to}\quad
x\le 1.
\]
Write \(g(x)=x-1\le 0\). The Lagrangian is
\[
\mathcal L(x,\lambda)=(x-2)^2+\lambda(x-1).
\]
KKT gives
\[
2(x-2)+\lambda=0,
\]
\[
x-1\le 0,\qquad \lambda\ge 0,\qquad \lambda(x-1)=0.
\]
The unconstrained minimizer \(x=2\) is infeasible, so the boundary should be active. Set \(x=1\). Stationarity gives
\[
2(1-2)+\lambda=0,
\]
so \(\lambda=2\). All KKT conditions hold, so \(x^\star=1\). The multiplier is positive because the boundary prevents the objective from moving toward its unconstrained minimizer.

\paragraph{Computational neuroscience example.}
Firing rates, probabilities, conductances, and variances are often constrained to be nonnegative. If a fitted parameter lands at zero, the usual gradient need not vanish. KKT conditions say that the gradient can be balanced by the nonnegativity constraint. For example, if a model estimates a nonnegative coupling strength as zero, the active constraint may explain why the optimum sits on the boundary.

\paragraph{AI and machine-learning example.}
Support vector machines are classic KKT examples. Inequality constraints enforce classification margins, and Lagrange multipliers become nonzero only for support vectors, the training points that lie on or violate the margin. In constrained fairness or resource-aware learning, KKT multipliers can indicate which constraints are binding at the optimum.

\paragraph{Common misconceptions and failure modes.}
\begin{itemize}[leftmargin=*]
\item Inactive inequalities have zero multipliers; they do not affect stationarity.
\item Positive multipliers are allowed only for constraints written in the form \(g_i(\theta)\le 0\). Changing the sign of \(g_i\) changes the sign convention.
\item KKT conditions are necessary under constraint qualifications; they are sufficient for global optimality mainly in convex settings.
\item Complementary slackness is a product condition, not a statement that both factors are zero.
\item Boundary optima are common in constrained statistical and neural models.
\end{itemize}

\paragraph{Exercises.}
\begin{enumerate}[label=\textbf{22.2.\arabic*.}, leftmargin=*]
\item \textbf{Mechanical.} For \(\min_x (x+1)^2\) subject to \(x\ge 0\), rewrite the inequality in \(g(x)\le 0\) form and write the Lagrangian.
\item \textbf{Proof or derivation.} Explain why complementary slackness implies that inactive inequality constraints have zero multipliers.
\item \textbf{Numerical/computational.} Solve \(\min_x (x+1)^2\) subject to \(x\ge 0\) using KKT conditions.
\item \textbf{Interpretation.} In a nonnegative firing-rate model, what does it mean if the optimal rate parameter is exactly zero and the associated inequality multiplier is positive?
\item \textbf{Misconception/debugging.} A student writes a multiplier \(\lambda=-3\) for an inequality \(g(x)\le 0\). Explain why this violates dual feasibility and how sign conventions matter.
\end{enumerate}

\subsection{22.3 Duality}

\paragraph{Guiding question.}
How can constraints produce lower bounds, alternative optimization problems, and interpretable shadow prices?

\paragraph{Beginner-facing intuition.}
The primal problem is the original constrained problem. The dual problem is built from the Lagrangian by asking: for fixed multipliers, what is the best lower bound on the primal objective? Weak duality says every dual feasible point gives a bound. Strong duality says that, under special conditions, the best bound equals the primal optimum.

Dual variables often have an economic or scientific interpretation. They measure how much the optimum would improve or worsen if a constraint were relaxed. In machine learning, duality also changes the computational form of some problems, exposing kernels, support vectors, or variational lower bounds.

\paragraph{Formal definitions and notation.}
Consider the primal problem
\[
p^\star=
\inf_{\theta\in\mathbb R^d} f(\theta)
\quad\text{subject to}\quad
g_i(\theta)\le 0,\quad h_j(\theta)=0.
\]
The Lagrangian is
\[
\mathcal L(\theta,\lambda,\nu)
=f(\theta)+\sum_i\lambda_i g_i(\theta)+\sum_j\nu_j h_j(\theta),
\qquad
\lambda_i\ge 0.
\]
The dual function is
\[
\boxed{
q(\lambda,\nu)=\inf_{\theta\in\mathbb R^d}\mathcal L(\theta,\lambda,\nu).
}
\]
The dual problem is
\[
\boxed{
d^\star=\sup_{\lambda\ge 0,\nu} q(\lambda,\nu).
}
\]
The difference \(p^\star-d^\star\) is the duality gap. Weak duality always gives
\[
d^\star\le p^\star.
\]
Strong duality means \(d^\star=p^\star\). For convex problems with affine equalities and suitable regularity, such as Slater's condition for strict feasibility of inequalities, strong duality often holds.

\paragraph{Core mechanism: proof of weak duality.}
Let \(\theta\) be any primal feasible point, so \(g_i(\theta)\le 0\) and \(h_j(\theta)=0\). Let \(\lambda_i\ge 0\). Then
\[
\sum_i\lambda_i g_i(\theta)\le 0,
\qquad
\sum_j\nu_j h_j(\theta)=0.
\]
Therefore
\[
\mathcal L(\theta,\lambda,\nu)
=f(\theta)+\sum_i\lambda_i g_i(\theta)+\sum_j\nu_j h_j(\theta)
\le f(\theta).
\]
Since \(q(\lambda,\nu)\) is the infimum of \(\mathcal L\) over all \(\theta\),
\[
q(\lambda,\nu)\le \mathcal L(\theta,\lambda,\nu)\le f(\theta)
\]
for every feasible \(\theta\). Thus \(q(\lambda,\nu)\) is a lower bound on every feasible objective value, so it is a lower bound on \(p^\star\). Maximizing over dual feasible multipliers preserves the bound:
\[
d^\star\le p^\star.
\]

\paragraph{Concrete worked example.}
Solve the primal problem
\[
\min_x \frac12x^2
\quad\text{subject to}\quad
x\ge 1.
\]
Write the constraint as
\[
g(x)=1-x\le 0.
\]
The Lagrangian is
\[
\mathcal L(x,\lambda)=\frac12x^2+\lambda(1-x),
\qquad \lambda\ge 0.
\]
For fixed \(\lambda\), minimize over \(x\):
\[
\frac{d}{dx}\mathcal L=x-\lambda=0
\quad\Rightarrow\quad
x=\lambda.
\]
Substitute:
\[
q(\lambda)=\frac12\lambda^2+\lambda(1-\lambda)
=\lambda-\frac12\lambda^2.
\]
The dual problem is
\[
\max_{\lambda\ge 0}\ \lambda-\frac12\lambda^2.
\]
The maximum occurs at \(\lambda=1\), giving
\[
d^\star=\frac12.
\]
The primal optimum is \(x^\star=1\), with value \(p^\star=1/2\), so strong duality holds here.

\paragraph{Computational neuroscience example.}
Suppose a neural population model minimizes decoding error subject to an average energy budget. The dual multiplier on the energy constraint measures the marginal value of energy: if the multiplier is large, relaxing the energy budget would substantially improve the objective. This interpretation is a property of the mathematical model, not a claim that a neural circuit explicitly computes a dual variable.

\paragraph{AI and machine-learning example.}
In support vector machines, the dual problem expresses the classifier in terms of inner products between training examples. Replacing inner products by kernels yields nonlinear decision boundaries without explicitly constructing high-dimensional features. In variational inference, a different kind of dual or variational viewpoint appears: a difficult log evidence is bounded below by an optimized functional, the ELBO. The theme is shared: construct a tractable bound whose tightness depends on an auxiliary object.

\paragraph{Common misconceptions and failure modes.}
\begin{itemize}[leftmargin=*]
\item Weak duality always gives a bound, but the bound may be loose.
\item Strong duality is not automatic for nonconvex problems.
\item The dual problem is not always easier than the primal; it is an alternative view.
\item A multiplier's shadow-price interpretation is local and model-dependent.
\item The sign of a dual bound depends on whether the primal is a minimization or maximization problem and on the constraint convention.
\end{itemize}

\paragraph{Exercises.}
\begin{enumerate}[label=\textbf{22.3.\arabic*.}, leftmargin=*]
\item \textbf{Mechanical.} Define the primal value \(p^\star\), dual function \(q\), dual value \(d^\star\), and duality gap.
\item \textbf{Proof or derivation.} Prove weak duality for a minimization problem with inequality constraints \(g_i(\theta)\le 0\).
\item \textbf{Numerical/computational.} Derive the dual function for \(\min_x \frac12(x-2)^2\) subject to \(x\le 0\), using \(g(x)=x\le 0\).
\item \textbf{Interpretation.} What would a large positive dual multiplier mean for an energy constraint in a neural decoder?
\item \textbf{Misconception/debugging.} A student says, ``The dual optimum is always equal to the primal optimum.'' State conditions under which this is true and give a reason it can fail.
\end{enumerate}

\subsection{22.4 Variational calculus}

\paragraph{Guiding question.}
How do we optimize over entire functions or distributions rather than finite-dimensional parameter vectors?

\paragraph{Beginner-facing intuition.}
Ordinary calculus asks how a scalar objective changes when a vector is nudged. Variational calculus asks how a functional changes when a whole function is nudged. The decision object might be a path \(y(t)\), a control signal \(u(t)\), a probability density \(q(z)\), or a neural field \(u(x)\). A small perturbation is not a vector step but a function \(\eta\) added to the candidate function.

The first variation plays the role of the gradient. At a stationary function, every allowed first-order perturbation has zero first-order effect. When the functional is an integral depending on a function and its derivative, this stationarity leads to the Euler-Lagrange equation.

\paragraph{Formal definitions and notation.}
A functional is a map from functions to numbers. A common form is
\[
J[y]=\int_a^b F(t,y(t),y'(t))\,dt,
\]
where \(y:[a,b]\to\mathbb R\) is differentiable and satisfies boundary conditions such as \(y(a)=A\), \(y(b)=B\). A variation is a perturbation
\[
y_\varepsilon(t)=y(t)+\varepsilon\eta(t),
\]
where \(\eta(a)=\eta(b)=0\) if endpoints are fixed.

The first variation is
\[
\delta J[y;\eta]
=\left.\frac{d}{d\varepsilon}J[y+\varepsilon\eta]\right|_{\varepsilon=0}.
\]
When it can be written as
\[
\delta J[y;\eta]=\int_a^b \frac{\delta J}{\delta y}(t)\eta(t)\,dt,
\]
the function \(\delta J/\delta y\) is the functional derivative. A stationary point satisfies
\[
\delta J[y;\eta]=0
\]
for every allowed variation \(\eta\).

\paragraph{Core mechanism: Euler-Lagrange equation.}
Let
\[
J[y]=\int_a^b F(t,y(t),y'(t))\,dt,
\]
with fixed endpoints. Compute the first variation:
\[
\delta J[y;\eta]
=\int_a^b
\left(
\frac{\partial F}{\partial y}\eta
+\frac{\partial F}{\partial y'}\eta'
\right)dt.
\]
Integrate the \(\eta'\) term by parts:
\[
\int_a^b \frac{\partial F}{\partial y'}\eta'\,dt
=
\left[\frac{\partial F}{\partial y'}\eta\right]_a^b
-\int_a^b
\frac{d}{dt}\left(\frac{\partial F}{\partial y'}\right)\eta\,dt.
\]
The boundary term vanishes because \(\eta(a)=\eta(b)=0\). Therefore
\[
\delta J[y;\eta]
=\int_a^b
\left[
\frac{\partial F}{\partial y}
-
\frac{d}{dt}\left(\frac{\partial F}{\partial y'}\right)
\right]\eta(t)\,dt.
\]
The fundamental lemma of variational calculus says that if
\[
\int_a^b r(t)\eta(t)\,dt=0
\]
for all smooth perturbations \(\eta\) vanishing at the endpoints, then \(r(t)=0\) on the interval, under standard regularity assumptions. Thus stationary functions satisfy
\[
\boxed{
\frac{\partial F}{\partial y}
-
\frac{d}{dt}\left(\frac{\partial F}{\partial y'}\right)=0.
}
\]
This is the Euler-Lagrange equation. It is a stationarity condition, not by itself a proof of a minimum.

\paragraph{Concrete worked examples.}
\emph{Shortest path.} For curves \(y(t)\) between fixed endpoints in the plane, arc length is
\[
J[y]=\int_a^b \sqrt{1+(y')^2}\,dt.
\]
Here \(F=\sqrt{1+(y')^2}\) does not depend on \(y\). The Euler-Lagrange equation gives
\[
0-\frac{d}{dt}\left(\frac{y'}{\sqrt{1+(y')^2}}\right)=0.
\]
Thus \(y'/\sqrt{1+(y')^2}\) is constant, so \(y'\) is constant. The stationary path is a straight line.

\emph{Smoothing functional.} Let
\[
J[y]=\int_a^b
\left[
\frac12(y(t)-s(t))^2+\frac{\lambda}{2}(y'(t))^2
\right]dt,
\]
where \(s(t)\) is a noisy signal and \(\lambda>0\). Then
\[
\frac{\partial F}{\partial y}=y-s,
\qquad
\frac{\partial F}{\partial y'}=\lambda y'.
\]
The Euler-Lagrange equation is
\[
(y-s)-\lambda y''=0.
\]
The fitted function balances closeness to data with smoothness.

\paragraph{Variational inference mechanism.}
Let \(x\) be observed data, \(z\) a latent variable, and \(q(z)\) an approximate posterior density constrained by
\[
q(z)\ge 0,\qquad \int q(z)\,dz=1.
\]
For the stationarity derivation below, assume an interior problem on a fixed support \(S\): \(p(x,z)>0\) is integrable on \(S\), \(q(z)>0\) for \(z\in S\), \(q(z)=0\) outside \(S\), and variations \(q_\varepsilon=q+\varepsilon\eta\) are supported on \(S\) and small enough that \(q_\varepsilon\) remains nonnegative. This assumption keeps \(\log p(x,z)\) finite and makes the derivative of \(q\log q\) well defined pointwise. If the optimum lies on the boundary with \(q(z)=0\) on part of the allowed support, additional inequality conditions are needed.

The evidence lower bound, or ELBO, is
\[
\mathcal E[q]
=\int q(z)\log p(x,z)\,dz-\int q(z)\log q(z)\,dz.
\]
Equivalently,
\[
\log p(x)=\mathcal E[q]+D_{\mathrm{KL}}(q(z)\|p(z\mid x)).
\]
Since KL divergence is nonnegative, \(\mathcal E[q]\le \log p(x)\), with equality when \(q(z)=p(z\mid x)\).

The functional stationarity calculation also shows this. Enforce normalization with a multiplier \(\lambda\):
\[
\mathcal A[q]
=\mathcal E[q]+\lambda\left(\int q(z)\,dz-1\right).
\]
For a variation \(q_\varepsilon=q+\varepsilon\eta\) with \(\int\eta(z)\,dz=0\) or with the multiplier included, the first variation is
\[
\delta\mathcal A[q;\eta]
=\int
\left[
\log p(x,z)-(\log q(z)+1)+\lambda
\right]\eta(z)\,dz.
\]
Stationarity for all \(\eta\) gives
\[
\log q(z)=\log p(x,z)+\lambda-1.
\]
Thus
\[
q(z)\propto p(x,z),
\]
and normalization gives
\[
q(z)=p(z\mid x).
\]
Variational inference restricts \(q\) to a tractable family, so it generally finds the best approximation within that family rather than the exact posterior.

\paragraph{Computational neuroscience example.}
Neural field models may describe a spatial activity profile \(u(x,t)\). An energy functional can penalize deviations from preferred activity, spatial roughness, and interaction terms. A stationary bump state can be studied by setting a functional derivative to zero. This is a mathematical model of population activity patterns, not a claim that cortex literally solves the Euler-Lagrange equation at every moment.

Optimal control in motor neuroscience uses function-valued decisions too. A control \(u(t)\) may minimize
\[
J[u]=\int_0^T \left(\|x(t)-x_{\mathrm{target}}\|^2+\rho\|u(t)\|^2\right)dt
\]
subject to dynamics \(\dot x=f(x,u,t)\). Multipliers or adjoint variables enforce the dynamics, producing necessary conditions for optimal trajectories. Chapter 30 returns to this theme.

\paragraph{AI and machine-learning example.}
ELBO optimization in variational autoencoders, energy-based learning, score matching, diffusion models, and entropy-regularized reinforcement learning all use variational ideas. The decision object may be a distribution \(q\), a path distribution, a score function, or a policy. The common mathematical pattern is: define a functional, compute a first variation or gradient, enforce normalization or dynamics with constraints, then optimize within a tractable parameterization.

\paragraph{Common misconceptions and failure modes.}
\begin{itemize}[leftmargin=*]
\item A functional derivative is not an ordinary derivative with respect to a scalar; it describes sensitivity to perturbation functions.
\item Euler-Lagrange equations give stationary candidates, not automatic minima.
\item Boundary conditions are part of the problem; changing them changes the stationarity equations.
\item In variational inference, the unrestricted optimum is the exact posterior, but practical approximations restrict \(q\).
\item Free energy and ELBO language should not be used as vague biological explanation without specifying the variables, functional, constraints, and update rule.
\end{itemize}

\paragraph{Exercises.}
\begin{enumerate}[label=\textbf{22.4.\arabic*.}, leftmargin=*]
\item \textbf{Mechanical.} For \(J[y]=\int_0^1 \frac12(y')^2\,dt\) with fixed endpoints, identify \(F(t,y,y')\), \(\partial F/\partial y\), and \(\partial F/\partial y'\).
\item \textbf{Proof or derivation.} Derive the Euler-Lagrange equation for \(J[y]=\int_a^b F(t,y,y')\,dt\), including the integration-by-parts step and the role of fixed endpoints.
\item \textbf{Numerical/computational.} Discretize \(J[y]=\int_0^1 \frac12(y-s)^2+\frac{\lambda}{2}(y')^2\,dt\) on three interior grid points. Write the finite-dimensional quadratic objective using finite differences.
\item \textbf{Interpretation.} In the ELBO identity \(\log p(x)=\mathcal E[q]+D_{\mathrm{KL}}(q\|p(\cdot\mid x))\), why does maximizing \(\mathcal E[q]\) over an unrestricted set of densities recover the exact posterior?
\item \textbf{Misconception/debugging.} A student says, ``The brain minimizes free energy,'' without specifying a model. Rewrite the claim in a mathematically meaningful way by naming a distribution, functional, variables, and update process.
\end{enumerate}

\paragraph{Closing bridge.}
Constrained optimization and variational calculus turn geometry into equations: multipliers balance forbidden directions, KKT conditions handle active boundaries, duality builds bounds, and functional derivatives optimize over paths and distributions. The next chapters move from mathematical optimization to computation: automatic differentiation, numerical linear algebra, and dynamical systems implement these ideas at scale.

\clearpage
